\section{Stability Validation}\label{sec:stability}

We performed three categories of validation to confirm that the test results are reliable and the underlying code is correct.

\subsection{Consistency Runs}

The full CPU test suite was executed 3 consecutive times in the same container, with identical configuration:

\begin{table}[h]
\centering
\caption{Consistency run results}
\label{tab:consistency}
\begin{tabular}{lccc}
\toprule
Run & Passed & Failed & Skipped (GPU) \\
\midrule
1/3 & 7 & 0 & 3 \\
2/3 & 7 & 0 & 3 \\
3/3 & 7 & 0 & 3 \\
\bottomrule
\end{tabular}
\end{table}

No non-determinism, port conflicts, or state leakage between runs was observed.
The low benchmark variance (Table~\ref{tab:benchmarks}) further supports measurement stability.

\subsection{Memory Sanitizer Analysis (ASAN/UBSAN)}

The BLB codebase was rebuilt from source with AddressSanitizer and UndefinedBehaviorSanitizer instrumentation:
\begin{lstlisting}[language=bash]
cmake -DCMAKE_CXX_FLAGS="-fsanitize=address,undefined
  -fno-omit-frame-pointer"
  -DCMAKE_EXE_LINKER_FLAGS="-fsanitize=address,undefined"
\end{lstlisting}

Runtime options: \texttt{ASAN\_OPTIONS=detect\_leaks=0} (SEAL's static allocations produce false-positive leak reports) and \texttt{UBSAN\_OPTIONS=print\_stacktrace=1}.

\begin{table}[h]
\centering
\caption{ASAN/UBSAN results}
\label{tab:asan}
\begin{tabular}{lcc}
\toprule
Test & Exit Code & Sanitizer Errors \\
\midrule
\texttt{test\_ntt} & 0 & None \\
\texttt{test\_linear\_operator} & server=0, client=0 & None \\
\texttt{test\_fixpoint} & server=0, client=0 & None \\
\texttt{test\_matmul} & server=0, client=0 & None \\
\bottomrule
\end{tabular}
\end{table}

These four tests cover all major code paths: HEXL-accelerated NTT, SEAL CKKS operations, emp-ot IKNP extension, and HE matrix operations.
No heap buffer overflows, use-after-free, double-free, signed integer overflow, or null pointer dereferences were detected.

\subsection{Numerical Cross-Validation}

To verify that the MPC protocols produce cryptographically correct results (not merely non-crashing), we cross-validated C++ test output against Python reference implementations.

\paragraph{Fixed-point comparison (\texttt{test\_fixpoint}).}
The C++ test executes \texttt{less\_than\_constant(x, 3, bitwidth=16)} on 32 random values using the IKNP OT protocol, with results XOR-shared between the two parties and reconstructed on BOB's side.
A Python reference validates the logic for signed 16-bit two's complement representation:
\begin{itemize}[nosep]
  \item Values in $[0, 32767]$ are positive; values in $[32768, 65535]$ represent negative numbers.
  \item \texttt{less\_than\_constant(x, 3)} returns 1 iff $\text{signed}(x) < 3$.
  \item Edge cases: $0 < 3$ (true), $3 < 3$ (false), $65535 = -1 < 3$ (true), $32768 = -32768 < 3$ (true).
\end{itemize}
All 32 cases matched between C++ and Python.

\paragraph{HE element-wise multiply (\texttt{test\_linear\_operator}).}
The test computes $x^2$ and $x \times y$ element-wise on 8192 CKKS slots with signed inputs in $[-5, 5]$.
Both operations pass: reconstructed ciphertext values match expected plaintext values.
